% ============================================================
% Part 1 — Medium Worked Problems
% Topics: sine rule, cosine rule, area formula, inverse trig,
%         compound angle identities, general triangles
% ============================================================

% ------------------------------------------------------------------
\begin{problem}[The Geometry of a 10 cm Cube]
A solid cube $ABCDEFGH$ has side length $10$ cm. The base is the square
$ABCD$, the top face is $EFGH$, and $E$ is directly above $A$.

\begin{center}
\begin{tikzpicture}[scale=0.72, line join=round, line cap=round]
  \coordinate (A) at (0,0);
  \coordinate (B) at (4,0);
  \coordinate (C) at (5.45,1.35);
  \coordinate (D) at (1.45,1.35);
  \coordinate (E) at (0,4);
  \coordinate (F) at (4,4);
  \coordinate (G) at (5.45,5.35);
  \coordinate (H) at (1.45,5.35);
  \coordinate (M) at (3.45,5.35);

  \fill[bookpurple!5] (A)--(B)--(C)--(D)--cycle;
  \fill[bookred!5] (E)--(F)--(G)--(H)--cycle;
  \draw[dashed, gray] (A)--(D)--(C) (D)--(H);
  \draw[thick] (A)--(B)--(C)--(G)--(F)--(B);
  \draw[thick] (F)--(E)--(H)--(G);
  \draw[thick] (A)--(E);
  \draw[bookred, dashed, thick] (A)--(C);
  \draw[blue!70!black, very thick] (A)--(G);
  \draw[orange!90!black, very thick] (A)--(M);
  \fill (M) circle (2pt);

  \node[below left] at (A) {$A$};
  \node[below] at (B) {$B$};
  \node[right] at (C) {$C$};
  \node[left] at (D) {$D$};
  \node[left] at (E) {$E$};
  \node[above] at (F) {$F$};
  \node[above right] at (G) {$G$};
  \node[above left] at (H) {$H$};
  \node[above] at (M) {$M$};
  \node[below] at ($(A)!0.5!(B)$) {$10$ cm};
\end{tikzpicture}
\end{center}

\begin{enumerate}[(i)]
  \item By considering $\triangle ABC$, calculate the exact length of the face
        diagonal $AC$.
  \item The space diagonal $AG$ passes through the centre of the cube. By
        considering $\triangle ACG$, show that $AG = 10\sqrt{3}$ cm.
  \item Find $\angle GAC$, the angle the space diagonal makes with the base, to
        the nearest degree.
  \item Let $M$ be the midpoint of $GH$. Calculate $AM$ and find
        $\angle MAC$ to the nearest tenth of a degree.
\end{enumerate}
\end{problem}

\begin{hintbox}
Use two-dimensional slices of the cube. First work in the base square, then use
the right-angled triangle containing $A$, $C$, and $G$. For the midpoint $M$,
coordinates such as $A=(0,0,0)$ are the cleanest way to find the needed lengths.
\end{hintbox}

\begin{solution}
\begin{enumerate}[(i)]
  \item In the square base,
  \[
  AC^2 = 10^2+10^2=200,
  \qquad AC=10\sqrt{2}\text{ cm}.
  \]
  \item Since $CG=10$ cm and $CG$ is perpendicular to the base,
  \[
  AG^2=AC^2+CG^2=(10\sqrt2)^2+10^2=300,
  \]
  so $AG=10\sqrt3$ cm.
  \item In $\triangle ACG$,
  \[
  \tan\angle GAC=\frac{CG}{AC}=\frac{10}{10\sqrt2}=\frac1{\sqrt2}.
  \]
  Hence $\angle GAC\approx 35^\circ$.
  \item Put $A=(0,0,0)$, $C=(10,10,0)$, $G=(10,10,10)$ and
  $H=(0,10,10)$. Then $M=(5,10,10)$, so
  \[
  AM=\sqrt{5^2+10^2+10^2}=15\text{ cm}.
  \]
  Also
  \[
  CM=\sqrt{(10-5)^2+(10-10)^2+(0-10)^2}=5\sqrt5.
  \]
  Applying the cosine rule in $\triangle AMC$,
  \[
  \cos\angle MAC
  =\frac{AM^2+AC^2-CM^2}{2(AM)(AC)}
  =\frac{225+200-125}{2(15)(10\sqrt2)}
  =\frac1{\sqrt2}.
  \]
  Therefore $\angle MAC=45.0^\circ$.
\end{enumerate}
\end{solution}

\begin{takeaways}
\begin{itemize}[leftmargin=*]
  \item In 3D trigonometry, find the two-dimensional triangle containing the
        length or angle you need.
  \item Coordinates are especially useful when midpoints or non-face diagonals
        are involved.
  \item A cuboid with dimensions $a$, $b$, and $c$ has space diagonal
        $\sqrt{a^2+b^2+c^2}$.
\end{itemize}
\end{takeaways}
% ------------------------------------------------------------------

\begin{problem}[Domain of Trigonometric Crossover Compositions]
Consider
\[
f(x)=\arcsin(2x-3),
\qquad
g(x)=\arccos\!\left(\frac{x}{2}\right).
\]
\begin{enumerate}[(i)]
  \item Determine the domains of $f$ and $g$ separately.
  \item Let $h(x)=g(f(x))$. Find the largest possible domain of $h$.
  \item Determine the range of $h$.
\end{enumerate}
\end{problem}

\begin{hintbox}
For a composite function $g(f(x))$, the output of $f$ must lie in the domain of
$g$. Use the principal range
$\arcsin u\in[-\pi/2,\pi/2]$ and remember that $\pi/2<2$.
\end{hintbox}

\begin{solution}
For $f$ to be defined,
\[
-1\le 2x-3\le 1
\quad\Longrightarrow\quad
1\le x\le 2.
\]
For $g$ to be defined,
\[
-1\le \frac{x}{2}\le 1
\quad\Longrightarrow\quad
-2\le x\le 2.
\]

Now $f$ maps $[1,2]$ onto $[-\pi/2,\pi/2]$. Since
\[
\left[-\frac{\pi}{2},\frac{\pi}{2}\right]\subset [-2,2],
\]
every output of $f$ is a valid input for $g$. Hence the domain of $h$ is
$[1,2]$.

As $x$ runs from $1$ to $2$, $f(x)$ runs from $-\pi/2$ to $\pi/2$. Since
$g(u)=\arccos(u/2)$ is decreasing in $u$,
\[
\operatorname{range}(h)
=
\left[
\arccos\!\left(\frac{\pi}{4}\right),
\arccos\!\left(-\frac{\pi}{4}\right)
\right].
\]
\end{solution}

\begin{takeaways}
\begin{itemize}[leftmargin=*]
  \item The domain of $g(f(x))$ is controlled by both the domain of $f$ and the
        requirement that $f(x)$ lies in the domain of $g$.
  \item Inverse trigonometric functions have restricted principal ranges; those
        ranges matter in composite functions.
  \item Monotonicity helps prevent endpoint reversals when finding a range.
\end{itemize}
\end{takeaways}
% ------------------------------------------------------------------

\begin{problem}[Overlapping Shadows]
Two right-angled triangles share a common horizontal base of length $12$. The
smaller triangle has angle of elevation $\pi/4$, while the larger triangle has
angle of elevation $\pi/3$. Let $x$ be the vertical distance between the two
peaks. Show that
\[
x=12(\sqrt3-1).
\]

\begin{center}
\begin{tikzpicture}[scale=0.36, line cap=round, line join=round]
  \coordinate (A) at (0,0);
  \coordinate (B) at (12,0);
  \coordinate (P) at (12,12);
  \coordinate (Q) at (12,20.784);

  \draw[thick] (A)--(B) node[midway, below] {$12$};
  \draw[thick] (B)--(Q);
  \draw[thick, bookpurple] (A)--(P);
  \draw[thick, bookred] (A)--(Q);
  \draw ($(B)+(-0.8,0)$)--++(0,0.8)--++(0.8,0);

  \draw (2.3,0) arc (0:45:2.3);
  \node at (3.2,1.1) {$\pi/4$};
  \draw (4.0,0) arc (0:60:4.0);
  \node at (4.2,3.0) {$\pi/3$};

  \draw[<->, thick] (13.0,12)--(13.0,20.784);
  \node[right] at (13.0,16.4) {$x$};
  \node[below left] at (A) {$A$};
  \node[below right] at (B) {$B$};
  \node[right] at (P) {$P$};
  \node[right] at (Q) {$Q$};
\end{tikzpicture}
\end{center}
\end{problem}

\begin{hintbox}
Find the two vertical heights separately using $\tan\theta=\frac{\text{opposite}}{\text{adjacent}}$,
then subtract them.
\end{hintbox}

\begin{solution}
For the smaller triangle,
\[
\tan\frac{\pi}{4}=\frac{BP}{12}.
\]
Since $\tan(\pi/4)=1$,
\[
BP=12.
\]
For the larger triangle,
\[
\tan\frac{\pi}{3}=\frac{BQ}{12}.
\]
Since $\tan(\pi/3)=\sqrt3$,
\[
BQ=12\sqrt3.
\]
The vertical distance between the peaks is therefore
\[
x=BQ-BP=12\sqrt3-12=12(\sqrt3-1).
\]
\end{solution}

\begin{takeaways}
\begin{itemize}[leftmargin=*]
  \item Shared-base height problems often reduce to two tangent ratios.
  \item Keep exact values in surd form for ``show that'' questions.
  \item Trigonometric functions are not linear: increasing the angle does not
        increase the height by the same factor.
\end{itemize}
\end{takeaways}
% ------------------------------------------------------------------

\begin{problem}[Inscribed Triangles and the Extended Sine Rule]
A triangle is inscribed in a circle. One side has length $12$ and passes through
the centre of the circle. Another side has length $7$.

\begin{enumerate}[(i)]
  \item A point $P$ is moved along the major arc of the circle. If the angle at
        $P$ is denoted by $\alpha$, explain why $\alpha$ remains constant and
        find $\cos\alpha$.
  \item Now suppose the triangle is changed so that its side lengths are
        $7$, $8$, and $10$. Find the angle $\theta$ opposite the side of length
        $8$, and determine the radius of its circumcircle.
\end{enumerate}

\begin{center}
\begin{tikzpicture}[scale=0.62, line cap=round, line join=round]
  \coordinate (O) at (0,0);
  \coordinate (A) at (-4,0);
  \coordinate (B) at (4,0);
  \coordinate (P) at (0,4);
  \draw[thick] (O) circle (4);
  \draw[thick, bookpurple] (A)--(B)--(P)--cycle;
  \draw[dashed] (O)--(A) (O)--(B);
  \fill (O) circle (2pt);
  \draw ($(P)+(0,-0.45)$) arc (-90:-135:0.45);
  \draw ($(P)+(0,-0.45)$) arc (-90:-45:0.45);
  \node[below] at (O) {$O$};
  \node[left] at (A) {$A$};
  \node[right] at (B) {$B$};
  \node[above] at (P) {$P$};
  \node[below] at ($(A)!0.5!(B)$) {$12$};
  \node[right] at ($(B)!0.5!(P)$) {$7$};
  \node at (0,3.25) {$\alpha$};

  \begin{scope}[xshift=10.5cm, scale=0.72]
    \coordinate (C) at (0,0);
    \coordinate (D) at (10,0);
    \coordinate (E) at (6.25,{5*sqrt(2)});
    \draw[thick, bookred] (C)--(D)--(E)--cycle;
    \node[below] at ($(C)!0.5!(D)$) {$10$};
    \node[right] at ($(D)!0.5!(E)$) {$8$};
    \node[left] at ($(E)!0.5!(C)$) {$7$};
    \draw (C) ++(0.9,0) arc (0:48.6:0.9);
    \node at (1.25,0.55) {$\theta$};
  \end{scope}
\end{tikzpicture}
\end{center}
\end{problem}

\begin{hintbox}
Part (i) uses the angle in a semicircle theorem. Part (ii) uses the cosine rule
first, then the extended sine rule $\frac{a}{\sin A}=2R$.
\end{hintbox}

\begin{solution}
Since the side of length $12$ passes through the centre, it is a diameter.
The angle subtended by a diameter at the circumference is a right angle, so
\[
\alpha=\frac{\pi}{2},
\qquad
\cos\alpha=0.
\]
The angle is constant because angles subtended by the same chord in the same
segment are equal.

For the triangle with side lengths $7$, $8$, and $10$, let $\theta$ be opposite
the side of length $8$. By the cosine rule,
\[
8^2=7^2+10^2-2(7)(10)\cos\theta.
\]
Thus
\[
64=149-140\cos\theta,
\qquad
\cos\theta=\frac{17}{28}.
\]
So
\[
\theta=\arccos\left(\frac{17}{28}\right)\approx0.919\text{ radians}.
\]
Also,
\[
\sin\theta
=\sqrt{1-\left(\frac{17}{28}\right)^2}
=\frac{\sqrt{495}}{28}
=\frac{3\sqrt{55}}{28}.
\]
Using the extended sine rule,
\[
2R=\frac{8}{\sin\theta}
=\frac{8}{3\sqrt{55}/28}
=\frac{224}{3\sqrt{55}},
\]
so
\[
R=\frac{112}{3\sqrt{55}}\approx5.03.
\]
\end{solution}

\begin{takeaways}
\begin{itemize}[leftmargin=*]
  \item A diameter subtends a right angle at the circumference.
  \item The cosine rule handles non-right-angled triangles when three sides are
        known.
  \item The extended sine rule links a triangle's side-angle data to its
        circumradius.
\end{itemize}
\end{takeaways}
